\documentclass[10pt,a4paper]{amsart}
\usepackage[margin=19mm]{geometry}
\usepackage{amsmath,amssymb,mathtools}
\usepackage[scr=rsfs,cal=euler]{mathalfa}
\usepackage{xfrac}
\usepackage[unicode,hidelinks,bookmarks=false]{hyperref}
\newcommand{\ie}{i.e.\ }
\newcommand{\cf}{cf.\ }
\newcommand{\resp}{respectively\ }
\newcommand{\Subsection}{\subsection}
\providecommand{\nobreakdash}{\nobreak\hbox{-}}
\setlength{\emergencystretch}{2em}
\multlinegap0pt
\usepackage{amssymb}
\usepackage{tikz}
\newtheorem{theorem}{Theorem}
\newtheorem{lemma}{Lemma}
\newcommand{\np}{\operatorname{np}}
\title{Some results on Archdeacon's conjecture for rotation systems}
\author{Arahat Chikkatur, Ji Zeng}
\date{Source: arXiv:2609.11599v1 (CC BY 4.0)}
\begin{document}
\maketitle

\noindent\emph{Source and licence.} Some results on Archdeacon's conjecture for rotation systems. Version arXiv:2609.11599v1 (CC BY 4.0), distributed by arXiv under the Creative Commons Attribution 4.0 International licence, http://creativecommons.org/licenses/by/4.0/, as recorded in the arXiv licence metadata for this exact version and shown on the abstract page https://arxiv.org/abs/2609.11599v1. Full licence notice: \texttt{licenses/arxiv-2609.11599-CC-BY-4.0.txt}.\par\smallskip

\noindent\textit{Editorial scope.} Counted unit: the article's theorem thm:antipodal (source0.tex lines 213--219). The article's complete original proof of it is delivered with it (source0.tex lines 731--812, one \texttt{proof} environment of 82 lines, which the article places away from the statement). No mathematics is written, repaired or re-derived: the statement and the proof are the article's own lines, in the article's own notation, and no retained line is substituted or re-flowed.  Retained uncounted as the source-local context the proof needs: lines 59--65; lines 307--320.  Nothing was omitted from any retained span, so the package carries no omission record.  No retained line contains a citation, so the wrapper carries no bibliography and no bibliography entry.  No retained line contains a figure environment, an image-inclusion command or drawing code, so no image is delivered and the package declares no asset dependency.  Editorial formatting only, in addition: an AMS \texttt{amsart} preamble that restates the article's own declarations and macros used by the retained lines, namely the environments \texttt{theorem}, \texttt{lemma} on separate counters, together with the standard packages those lines need.  The retained environments print in this document as Theorem 1, Lemma 1 in order of appearance, whereas the article prints the same items with the numbering of its own full document.  The complete native source of the article as distributed in the arXiv e-print for this exact version is bundled unchanged as \texttt{sources/210036/source0.tex}.
\begin{equation}\label{eq:hill-number}
H(n)=\frac14
\left\lfloor\frac n2\right\rfloor
\left\lfloor\frac{n-1}2\right\rfloor
\left\lfloor\frac{n-2}2\right\rfloor
\left\lfloor\frac{n-3}2\right\rfloor
\end{equation}

\begin{lemma}\label{planar}
Let \(R\) be a rotation system on a four-element set, and let
\((z_0,z_1,z_2,z_3)\) be any ordering of its ground set.  Then \(R\) is
planar if and only if
\begin{equation}\label{eq:planar-orientations}
\bigl(\varepsilon_{z_0}(z_1,z_2,z_3),\;
      \varepsilon_{z_1}(z_0,z_2,z_3),\;
      \varepsilon_{z_2}(z_0,z_1,z_3),\;
      \varepsilon_{z_3}(z_0,z_1,z_2)\bigr)
\in\bigl\{(0,1,0,1),\,(1,0,1,0)\bigr\}.
\end{equation}
In particular, whether this alternating pattern occurs does not depend
on which ordering of the ground set is used.
\end{lemma}

\begin{theorem}\label{thm:antipodal}
If \(\{u,v\}\) is an antipodal pair in a rotation system \(R\) on
\(n\) elements, then
\[
\np(R)\geq \np(R\setminus\{u,v\})+H(n)-H(n-2).
\]
\end{theorem}

\begin{proof}[Proof of Theorem~\ref{thm:antipodal}]
Let \(\{u,v\}\) be the given antipodal pair, and
\(U\) be the set of elements not equal to \(u\) or \(v\). Write \(m=|U|=n-2\). We identify \(U\)
with \([m]\) so that
\[
p_u=(v,1,\ldots,m)
\quad\text{and}\quad
p_v=(u,m,\ldots,1).
\]
Introduce an auxiliary symbol \(\star\), and let \(c_u^v\) be obtained
from \(p_u\) by replacing \(v\) with \(\star\), and \(c_v^u\) from
\(p_v\) by replacing \(u\) with \(\star\). By our hypothesis, the two cyclic orders \(c_u^v\) and \(c_v^u\) are reverses of each other.

We call a triple \(T\in\binom{U\cup\{\star\}}3\) \emph{good} if
\(T\cup\{u\}\) and \(T\cup\{v\}\) are planar in case
\(\star\notin T\), and if \(\{u,v\}\cup(T\setminus\{\star\})\) is
planar in case \(\star\in T\). For distinct \(a,b\in U\), deleting \(a\) and \(b\) from \(c_u^v\) gives us two intervals of elements; let \(I_{a,b}\) be the one containing \(\star\) and
\(J_{a,b}\) the other. Note that since \(c_v^u\) is the reverse of \(c_u^v\), deleting \(a\) and \(b\) from \(c_v^u\) gives us the same two intervals. An element \(x \in (U\cup\{\star\}) \setminus \{a,b\}\) is said to be a \emph{good complement} of \(\{a,b\}\) if the triple \(\{a,b,x\}\) is good.

\begin{lemma}\label{lem:good-interval}
For distinct \(a,b\in U\), at most one of the intervals \(I_{a,b}\) and \(J_{a,b}\) contains a good complement of \(\{a,b\}\).
\end{lemma}

\begin{proof}

Suppose on the contrary that \(x \in I_{a,b}\) and \(y \in J_{a,b}\) are two good complements of \(\{a,b\}\). We first deal with the case \(x \neq \star\). By our definition of \(I_{a,b}\) and \(J_{a,b}\), we have \(\varepsilon_u(a,b,x) = 1 - \varepsilon_u(a,b,y)\). Without loss of generality, we assume \(\varepsilon_u(a,b,x) = 0\) and \(\varepsilon_u(a,b,y) = 1\). Note that we also have \(\varepsilon_v(a,b,x) = 1 -\varepsilon_u(a,b,x)\) and similarly for \(y\). Because \(\{u,x,a,b\}\) and \(\{v,x,a,b\}\) are planar as \(x\) is a good complement, by Lemma~\ref{planar} applied to \((u,a,b,x)\) and \((v,a,b,x)\), we have \(\varepsilon_a(u,b,x)=1\) and \(\varepsilon_a(v,b,x)=0\). Similarly, because \(y\) is a good complement, we have \(\varepsilon_a(u,b,y)=0\) and \(\varepsilon_a(v,b,y)=1\).

Now consider the sequence of entries we encounter as we read the cyclic order \(p_a\) starting after and ending before \(b\): \(\varepsilon_a(u,b,x)=1\) means we encounter \(u\) before \(x\), \(\varepsilon_a(v,b,x)=0\) means we encounter \(x\) before \(v\), and together they imply that we encounter \(u\) before \(v\);  \(\varepsilon_a(u,b,y)=0\) means we encounter \(y\) before \(u\), \(\varepsilon_a(v,b,y)=1\) means we encounter \(v\) before \(y\), and together they imply that we encounter \(v\) before \(u\); hence a contradiction is reached.

It remains to deal with the case \(x=\star\), in which
\(\{u,v,a,b\}\) is planar.  Since \(\star\) takes the place of \(v\) in
\(c_u^v\), our definition of \(I_{a,b}\) and \(J_{a,b}\) gives
\(\varepsilon_u(a,b,v)=1-\varepsilon_u(a,b,y)\). Without loss of generality, we assume \(\varepsilon_u(a,b,v)=0\) and \(\varepsilon_u(a,b,y)=1\). Just as the previous case, we apply Lemma~\ref{planar} to \((u,a,b,v)\), \((u,a,b,y)\), and \((v,a,b,y)\), we can get \(\varepsilon_a(u,b,v)=1\), \(\varepsilon_a(u,b,y)=0\), and \(\varepsilon_a(v,b,y)=1\). Now we read the entries in \(p_a\) with \(b\) being the cut: \(\varepsilon_a(u,b,v)=1\) means we encounter \(u\) before \(v\);
\(\varepsilon_a(v,b,y)=1\) means we encounter \(v\) before \(y\), \(\varepsilon_a(u,b,y)=0\) means we encounter \(y\) before \(u\), and together they imply that we encounter \(v\) before \(u\); hence a contradiction is reached.
\end{proof}

Now we define a tournament \(\mathcal T\), i.e. a complete directed graph, on \([m]\) by setting, for
\(a<b\), \(a\to b\) if every good complement of \(\{a,b\}\) lies in
\(I_{a,b}\), and \(b\to a\) otherwise; by
Lemma~\ref{lem:good-interval} the second case means that every good
complement lies in \(J_{a,b}\).

Since \(\star \in I_{a,b}\) by definition, each edge
\(b\to a\) with \(a<b\) means that \(\star\) is not a good complement
of \(\{a,b\}\), and hence that \(\{u,v,a,b\}\) is non-planar. Now let \(a<b<c\) and suppose \(\{a,b,c\}\) is good.  Then we have
\(c\in I_{a,b}\), \(a\in I_{b,c}\), and \(b\in J_{a,c}\);
hence \(a\to b\to c\to a\).  So a good triple in \(U\) induces a
cyclic triangle in \(\mathcal{T}\), and therefore every non-cyclic triangle
\(\{a,b,c\}\) yields a non-planar four-subset among
\(\{u,a,b,c\}\) and \(\{v,a,b,c\}\). These four-subsets considered are pairwise
distinct: those coming from edges contain both \(u\) and \(v\) with the other two elements given by the edge, those
coming from non-cyclic triangles contain exactly one of \(u\) or \(v\) with the other three elements given by the triangle. Writing \(\beta(\mathcal T)\) for the number
of pairs \(a<b\) with \(b\to a\), and \(\tau(\mathcal T)\) for the
number of non-cyclic triangles, we have
\begin{equation}\label{eq:new-np-tournament}
\np(R)-\np(R[U])\geq \beta(\mathcal T)+\tau(\mathcal T).
\end{equation}

It remains to bound the right-hand side.  Let \(d_a\) be the
outdegree of \(a\) and \(\beta_a\) the number of its outneighbours among
\(\{1,\ldots,a-1\}\).  Every non-cyclic triangle has a unique source, and
any two outneighbours of an element form one with that element as
source, so \(\tau(\mathcal T)=\sum_a\binom{d_a}2\) and
\(\beta(\mathcal T)=\sum_a \beta_a\); since \(a\) has only \(m-a\)
successors, \(\beta_a\geq\max\{0,d_a-(m-a)\}\).  Defining
\(f_a(d)=\binom d2+\max\{0,d-(m-a)\}\) we obtain
\begin{equation}\label{eq:tournament-degree-bound}
\beta(\mathcal T)+\tau(\mathcal T)\geq\sum_{a=1}^m f_a(d_a),
\qquad\text{where }\ \sum_{a=1}^m d_a=\binom m2 .
\end{equation}

To bound the right-hand side of
\eqref{eq:tournament-degree-bound} from below, define
\(\Delta_a(d)=f_a(d)-f_a(d-1)\) for \(1\leq d\leq m-1\).
Let us arrange these numbers into an \(m\)-by-\((m-1)\) matrix \(\Delta\) whose \((a,d)\)-entry is \(\Delta_a(d)\).
Since \(f_a(0)=0\), we have \(f_a(d_a)=\sum_{d=1}^{d_a}\Delta_a(d)\). Then the right hand side of \eqref{eq:tournament-degree-bound} is a sum of \(\binom{m}{2}\) entries of the matrix \(\Delta\), no two from the same position. We can check by definition that \(\Delta_a(d) = d-1\) for \(d \leq m-a\) and \(\Delta_a(d) = d\) for \(d > m-a\). Hence, the \(a\)-th row of \(\Delta\) has entries \(0,1,\ldots,m-1\) with \(m-a\) removed, hence each integer \(0,1,\ldots,m-1\) appears exactly \(m-1\) times in \(\Delta\). Taking the \(\binom{m}{2}\) smallest entries from \(\Delta\) and summing them up, we get a lower bound \begin{equation}\label{eq:telescoping-lower-bound}
    \sum_{a=1}^m f_a(d_a) \geq \begin{cases}
       (2q-1)q(q-1)/2 & \text{if \(m=2q\),}\\
       q^3 & \text{if \(m=2q+1\).}
    \end{cases}
\end{equation} Using \eqref{eq:hill-number}, we can check that the right hand side of \eqref{eq:telescoping-lower-bound} is exactly \(H(m+2) - H(m)\). Since \(m = n-2\), combining this with \eqref{eq:new-np-tournament} and \eqref{eq:tournament-degree-bound} gives \(\np(R) \geq \np(R[U]) + H(n) - H(n-2)\), as required.
\end{proof}

\end{document}
